\chapter{Alcoholen: de OH-groep activeren}\label{ch:b1:alcohol-activation}

Alcoholen zijn overal: ze worden per ton gemaakt uit alkenen en door
fermentatie, en elke grignardadditie van het vorige hoofdstuk levert er
een op. Toch ondergaat een alcohol geen substitutie zoals een
halogeenalkaan: de OH-groep zou moeten vertrekken als het hydroxide-ion,
een \omterm{def:b1:predominant-reaction:strong-acid}{sterke base} en een zeer slechte \omterm{def:b1:electronic-effects:nucleophile}{vertrekkende groep}
(\cref{ch:b1:nucleophilic-substitution}). Chemici hebben drie
antwoorden. Verwijder in plaats daarvan het proton van de OH, en het
\omterm{def:b1:alcohol-activation:alkoxide}{alkoxide} wordt een krachtig \omterm{def:b1:electronic-effects:nucleophile}{nucleofiel}. Voeg een proton toe, en de
\omterm{def:b1:electronic-effects:nucleophile}{vertrekkende groep} wordt water. Of vervang het waterstofatoom door een
groep die de zuurstof deel maakt van een uitstekende \omterm{def:b1:electronic-effects:nucleophile}{vertrekkende groep},
zonder aan het koolstofatoom te raken. Dit hoofdstuk volgt de drie
wegen, en de keuzes van mechanisme en stereochemie die erbij horen.

\begin{recall}
SN1, SN2 en hun stereochemie (\cref{ch:b1:nucleophilic-substitution});
E1, E2 en de \omterm{prop:b1:elimination:zaitsev}{regel van Zaitsev} (\cref{ch:b1:elimination}); \omterm{def:b1:predominant-reaction:acidity-constant}{zuurconstanten}
en de methode van de \omterm{def:b1:predominant-reaction:predominant-reaction}{predominante reactie} (\cref{ch:b1:predominant-reaction});
het effect van alkylgroepen op de zuurheid van alcoholen
(\cref{exo:b1:electronic-effects:11}).
\end{recall}

\section{Alcoholen als zuren: alkoxiden}

\begin{definition}[Alkoxide]\label{def:b1:alcohol-activation:alkoxide}
Een \emph{alkoxide}\index{alkoxide} is de geconjugeerde base \ce{RO-} van
een alcohol \ce{ROH}: methoxide \ce{CH3O-}, ethoxide \ce{CH3CH2O-},
tert-butoxide \ce{(CH3)3CO-}. Alkoxiden zijn \omterm{def:b1:predominant-reaction:strong-acid}{sterke basen} en goede
\omterm{def:b1:electronic-effects:nucleophile}{nucleofielen}.
\end{definition}

\begin{proposition}[Zuurheid van alcoholen]\label{prop:b1:alcohol-activation:alcohol-acidity}
Alcoholen zijn zwakkere zuren dan water of ongeveer even zwak:
$\mathrm{p}K_a$ 15.5 (methanol), 15.9 (ethanol), 17.1 (propaan-2-ol),
19.2 (2-methylpropaan-2-ol), tegenover 14.0 voor het koppel
\ce{H2O}/\ce{OH-}. Hydroxide deprotoneert ze dus maar gedeeltelijk; een
volledige omzetting in het \omterm{def:b1:alcohol-activation:alkoxide}{alkoxide} vraagt een sterkere base of een
alkalimetaal.
\end{proposition}
% ledger: pka:methanol, pka:ethanol, pka:2-propanol, pka:tBuOH

\begin{proof}
De reactie \ce{ROH + OH- <=> RO- + H2O} heeft $K = 10^{14.0 -
\mathrm{p}K_a(\ce{ROH})}$, tussen $10^{-1.5}$ en $10^{-5}$: ongunstig.
Natriummetaal reduceert het proton, \ce{2ROH + 2Na -> 2RO- + 2Na+ + H2},
en het hydride-ion van natriumhydride neemt het onomkeerbaar op,
\ce{ROH + H- -> RO- + H2}: de waterstof ontsnapt, en de reactie is
volledig, wat haar evenwichtsconstante ook is.
\end{proof}

\begin{inthelab}[Natriumhydride]
Natriumhydride wordt verkocht als een grijs poeder, gedispergeerd in
minerale olie, die het tegen het vocht van de lucht beschermt. Het
reageert heftig met water en maakt daarbij brandbare waterstof vrij; het
wordt snel afgewogen, in porties onder een inert gas aan de droge alcohol
toegevoegd, en het opborrelen van waterstof toont hoe ver de deprotonering
gevorderd is. Overmaat hydride wordt aan het eind vernietigd door langzaam
een alcohol toe te voegen, nooit water.
\end{inthelab}

\section{De ethersynthese van Williamson}

\begin{definition}[Ethersynthese van Williamson]\label{def:b1:alcohol-activation:williamson}
De \emph{ethersynthese van Williamson}\index{ethersynthese van Williamson}
bereidt een ether $\mathrm{R{-}O{-}R'}$ door een SN2-reactie van een \omterm{def:b1:alcohol-activation:alkoxide}{alkoxide} \ce{RO-}
op een halogeenalkaan (of een \omterm{def:b1:alcohol-activation:sulfonate}{sulfonaatester}) $\mathrm{R'{-}X}$.
\end{definition}

\begin{omfigure}
\setchemfig{atom sep=2.1em}
\schemestart
\chemfig{CH_3CH_2-@{o}O^{-}}\,\chemfig{Na^{+}}
\+
\chemfig{@{c}CH_3-[@{ci}]@{i}I}
\arrow{->}
\chemfig{CH_3CH_2-O-CH_3}
\+
\chemfig{Na^{+}}\,\chemfig{I^{-}}
\schemestop

\medskip
\chemmove{\draw[->, thick, shorten <=2pt, shorten >=2pt](o) .. controls +(80:8mm) and +(100:8mm) .. (c);
          \draw[->, thick, shorten <=2pt, shorten >=2pt](ci) .. controls +(90:5mm) and +(90:5mm) .. (i);}

{\small Williamsonsynthese van methoxyethaan: het ethoxide valt het
koolstofatoom van joodmethaan aan aan de kant tegenover het jood (SN2).}
\end{omfigure}

\begin{method}[De twee partners kiezen]\label{met:b1:alcohol-activation:williamson-choice}
Een ether $\mathrm{R{-}O{-}R'}$ kan aan beide kanten van de zuurstof geknipt worden.
\begin{enumerate}
  \item Schrijf de twee mogelijke paren op: \ce{RO-} met $\mathrm{R'X}$, en
        $\mathrm{R'O^-}$ met \ce{RX}.
  \item Houd het paar waarin het halogenide methyl of primair is: SN2 heeft
        een onbelemmerd koolstofatoom nodig.
  \item Verwerp een paar met een tertiair halogenide: het \omterm{def:b1:alcohol-activation:alkoxide}{alkoxide}, een
        \omterm{def:b1:predominant-reaction:strong-acid}{sterke base}, zou in plaats daarvan elimineren (E2).
  \item Verwacht bij een secundair halogenide een mengsel van ether en
        alkeen.
\end{enumerate}
\end{method}

\begin{example}[Een asymmetrische ether]\label{ex:b1:alcohol-activation:williamson}
2-Methoxy-2-methylpropaan (methyl-tert-butylether): natrium-tert-butoxide
met joodmethaan geeft het zuiver; natriummethoxide met het tertiaire
bromide, 2-broom-2-methyl\-propaan, geeft 2-methylpropeen via E2. Omvangrijk
aan de kant van het \omterm{def:b1:alcohol-activation:alkoxide}{alkoxide}, klein aan de kant van het halogenide.
\end{example}

\section{De OH-groep activeren}

\begin{definition}[Sulfonaatesters]\label{def:b1:alcohol-activation:sulfonate}
Een \emph{sulfonaatester}\index{sulfonaatester} $\mathrm{R{-}O{-}SO_2{-}R'}$ ontstaat
uit een alcohol en een sulfonylchloride $\mathrm{R'SO_2Cl}$ in aanwezigheid van een
base (pyridine). Het \emph{tosylaat}\index{tosylaat}, \ce{R-OTs}, komt
van 4-methylbenzeensulfonylchloride (tosylchloride); het
\emph{mesylaat}\index{mesylaat}, \ce{R-OMs}, van methaansulfonylchloride.
Het sulfonaatanion $\mathrm{R'SO_3^-}$, de geconjugeerde base van een \omterm{def:b1:predominant-reaction:strong-acid}{sterk zuur} en
door resonantie over drie zuurstofatomen gestabiliseerd, is een
uitstekende \omterm{def:b1:electronic-effects:nucleophile}{vertrekkende groep}.
\end{definition}

\begin{omfigure}
\setchemfig{atom sep=2em}
\schemestart
\chemfig{R-O-H}
\+
\chemfig{Cl-S(=[2]O)(=[6]O)-[:0]*6(-=-(-CH_3)=-=)}
\arrow{->[pyridine]}
\chemfig{R-O-S(=[2]O)(=[6]O)-[:0]*6(-=-(-CH_3)=-=)}
\schemestop
\par\vspace{6pt}

{\small Tosylering van een alcohol. De O--H-binding wordt vervangen door
O--S; de C--O-binding van de alcohol blijft onaangeroerd. Pyridine neemt
het gevormde HCl op.}
\end{omfigure}

\begin{proposition}[Configuratie bij tosylering en substitutie]\label{prop:b1:alcohol-activation:tosylation-retention}
De tosylering van een alcohol op een stereogeen koolstofatoom behoudt zijn
\omterm{def:b1:stereochemistry-in-depth:isomers}{configuratie} (retentie); een daaropvolgende SN2-substitutie van het
\omterm{def:b1:alcohol-activation:sulfonate}{tosylaat} keert haar om. Samen vervangen de twee stappen OH door een
\omterm{def:b1:electronic-effects:nucleophile}{nucleofiel} met inversie.
\end{proposition}

\begin{proof}
Bij de tosylering veranderen de bindingen O--H en S--Cl: de zuurstof van
de alcohol valt het zwavelatoom aan en zijn waterstofatoom gaat naar de
base. Er wordt geen enkele binding met het stereogene koolstofatoom
verbroken, dus de schikking eromheen blijft dezelfde. In de SN2-stap
breekt de C--O-binding terwijl het \omterm{def:b1:electronic-effects:nucleophile}{nucleofiel} aan de tegenovergestelde
kant bindt (\cref{prop:b1:nucleophilic-substitution:sn2-inversion}).
\end{proof}

\begin{omfigure}
\begin{tikzpicture}[font=\small, >={Stealth}]
  \node (a) at (0,0) {\chemfig{C(-[:150]H_3C)(<[:210]H)(<:[:250]C_2H_5)-OH}};
  \node at (0,-1.4) {\cip{S}-butaan-2-ol};
  \draw[->, thick] (1.3,0) -- (2.5,0) node[midway, above] {TsCl} node[midway, below, font=\scriptsize] {retentie};
  \node (b) at (3.9,0) {\chemfig{C(-[:150]H_3C)(<[:210]H)(<:[:250]C_2H_5)-OTs}};
  \node at (3.9,-1.4) {\cip{S}-tosylaat};
  \draw[->, thick] (5.3,0) -- (6.6,0) node[midway, above] {\ce{CN-}} node[midway, below, font=\scriptsize] {SN2, inversie};
  \node (c) at (8.3,0) {\chemfig{[:0]NC-C(-[:30]CH_3)(<[:-30]H)(<:[:-70]C_2H_5)}};
  \node at (8.3,-1.4) {\cip{R}-nitril};
\end{tikzpicture}

{\small Van \cip{S}-butaan-2-ol naar 2-methylbutaannitril: de tosylering
behoudt de \omterm{def:b1:stereochemistry-in-depth:isomers}{configuratie}, de SN2-stap met cyanide keert haar om. Netto:
OH vervangen door CN met inversie.}
\end{omfigure}

Hetzelfde doel kan in zuur milieu bereikt worden: de geprotoneerde
OH-groep vertrekt als water, \ce{R-OH2+ -> R+ + H2O} (SN1, E1) of onder
aanval (SN2). Zuur beperkt het \omterm{def:b1:electronic-effects:nucleophile}{nucleofiel} echter tot de deeltjes die het
overleven --- halogenide-ionen, water, alcoholen --- en zet de deur open
voor \omterm{def:b1:electronic-effects:intermediates}{carbokationen} en hun nevenreacties.

\section{Omzetting in halogeenalkanen}

\begin{proposition}[Alcoholen en waterstofhalogeniden]\label{prop:b1:alcohol-activation:hx-by-class}
Tertiaire alcoholen geven het halogeenalkaan met geconcentreerd zoutzuur
bij kamertemperatuur (SN1); primaire alcoholen hebben
broomwaterstofzuur of joodwaterstofzuur en verwarming nodig (SN2 op de
geprotoneerde alcohol); secundaire alcoholen reageren met tussenliggende
snelheden, vaak volgens beide mechanismen.
\end{proposition}

\begin{proof}
Protonering maakt van OH een \ce{-OH2+}-groep. Bij een tertiaire alcohol
vertrekt water gemakkelijk, wat een stabiel \omterm{def:b1:electronic-effects:intermediates}{carbokation} geeft dat door
het halogenide ingevangen wordt (SN1): zelfs het bescheiden \omterm{def:b1:electronic-effects:nucleophile}{nucleofiel}
\ce{Cl-} volstaat. Een primair \omterm{def:b1:electronic-effects:intermediates}{carbokation} is te instabiel; het
halogenide moet water van achteren wegduwen (SN2), wat een goed
\omterm{def:b1:electronic-effects:nucleophile}{nucleofiel} vraagt (\ce{Br-}, \ce{I-}, in protische media beter dan
\ce{Cl-}, \cref{prop:b1:electronic-effects:nucleophilicity-basicity}) en
warmte.
\end{proof}

\begin{example}[Een classificatietest]\label{ex:b1:alcohol-activation:lucas}
Een mengsel van geconcentreerd zoutzuur en zinkchloride (een \omterm{def:b1:lewis-resonance-vsepr:lewis-acid}{lewiszuur}
dat de OH helpt vertrekken) maakt een tertiaire alcohol meteen troebel ---
het onoplosbare chloride scheidt zich af ---, een secundaire alcohol
binnen enkele minuten, en een primaire alcohol bij kamertemperatuur
nauwelijks: de volgorde van stabiliteit van de \omterm{def:b1:electronic-effects:intermediates}{carbokationen}, zichtbaar
gemaakt in een reageerbuis.
\end{example}

\section{Dehydratatie en ethervorming in zuur milieu}

\begin{definition}[Dehydratatie van een alcohol]\label{def:b1:alcohol-activation:dehydration}
De \emph{dehydratatie}\index{dehydratatie} van een alcohol is de
eliminatie van water, gekatalyseerd door een \omterm{def:b1:predominant-reaction:strong-acid}{sterk zuur} (zwavelzuur of
fosforzuur) en warmte, waarbij een alkeen ontstaat: $\mathrm{R_2CH{-}CR'_2OH} \to \mathrm{R_2C{=}CR'_2} + \ce{H2O}$.
\end{definition}

\begin{omfigure}
\setchemfig{atom sep=2.1em}
\schemestart
\chemfig{C(-[2]CH_3)(-[:210]H_3C)(-[:-30]CH_2CH_3)-OH}
\arrow{->[\ce{H+}]}
\chemfig{C(-[2]CH_3)(-[:210]H_3C)(-[:-30]CH_2CH_3)-OH_2^{+}}
\arrow{->[\ce{-H2O}]}
\chemfig{C^{+}(-[2]CH_3)(-[:210]H_3C)(-[:-30]CH_2CH_3)}
\schemestop

\medskip
\schemestart
\chemfig{C^{+}(-[2]CH_3)(-[:210]H_3C)(-[:-30]CH_2CH_3)}
\arrow{->[\ce{-H+}][(hoofdproduct)]}
\chemfig{C(-[2]CH_3)(-[:210]H_3C)=[:-30]CHCH_3}
\schemestop

\medskip
{\small Zuurgekatalyseerde \omterm{def:b1:alcohol-activation:dehydration}{dehydratatie} van 2-methylbutaan-2-ol (E1):
protonering, verlies van water tot het tertiaire \omterm{def:b1:electronic-effects:intermediates}{carbokation}, verlies van
een proton. Het proton van de \ce{CH2} verwijderen geeft het
trigesubstitueerde 2-methylbut-2-een (Zaitsev, hoofdproduct); van een
methylgroep, 2-methylbut-1-een (nevenproduct).}
\end{omfigure}

\begin{proposition}[Ether of alkeen]\label{prop:b1:alcohol-activation:ether-vs-alkene}
Een primaire alcohol die met zwavelzuur verwarmd wordt, geeft bij matige
temperatuur vooral een symmetrische ether (het ene molecuul substitueert
het andere via SN2) en bij hogere temperatuur vooral een alkeen;
secundaire en tertiaire alcoholen geven alkenen.
\end{proposition}

\begin{proof}
Bij een primaire alcohol kan de geprotoneerde alcohol aangevallen worden
op het koolstofatoom door een tweede alcoholmolecuul (SN2, wat na verlies
van een proton \ce{R-O-R} geeft) of op een $\beta$-waterstofatoom (E2).
Eliminatie maakt meer moleculen dan ze verbruikt en wordt begunstigd
naarmate de temperatuur stijgt (\cref{prop:b1:elimination:sn-vs-e});
substitutie overheerst wanneer die lager is. Secundaire en tertiaire
alcoholen ioniseren tot \omterm{def:b1:electronic-effects:intermediates}{carbokationen}, die gemakkelijker een proton
verliezen dan dat ze door een zwak \omterm{def:b1:electronic-effects:nucleophile}{nucleofiel} ingevangen worden, des te
meer bij verwarming; het gevormde water wordt bovendien verwijderd door
het vluchtige alkeen af te destilleren.
\end{proof}

\section{Oefeningen}

\begin{exercise}[$\star$]\label{exo:b1:alcohol-activation:1}
Schrijf de reacties op van ethanol met natrium, met natriumhydride en met
natriumhydroxide; bereken de constante van de laatste ($\mathrm{p}K_a$ van
ethanol 15.9).
\end{exercise}

\begin{exercise}[$\star$]\label{exo:b1:alcohol-activation:2}
Geef de producten van: natriummethoxide met 1-broompropaan;
natriumethoxide met joodmethaan; natriumfenoxide met broomethaan.
\end{exercise}

\begin{exercise}[$\star$]\label{exo:b1:alcohol-activation:3}
Schrijf de tosylering van propaan-1-ol op en de reactie van het \omterm{def:b1:alcohol-activation:sulfonate}{tosylaat}
met natriumjodide. Wat is het voordeel tegenover de rechtstreekse reactie
van de alcohol met natriumjodide?
\end{exercise}

\begin{exercise}[$\star$]\label{exo:b1:alcohol-activation:4}
Geef de alkenen die ontstaan bij de \omterm{def:b1:alcohol-activation:dehydration}{dehydratatie} in zuur milieu van
butaan-2-ol en van 2-methylpentaan-2-ol, en het hoofdproduct.
\end{exercise}

\begin{exercise}[$\star\star$]\label{exo:b1:alcohol-activation:5}
Stel een williamsonsynthese voor van 1-ethoxy-2-methylpropaan en verklaar
je keuze van de partners. Waarom zou de andere keuze een alkeen geven?
\end{exercise}

\begin{exercise}[$\star\star$]\label{exo:b1:alcohol-activation:6}
\cip{R}-octaan-2-ol wordt omgezet in zijn \omterm{def:b1:alcohol-activation:sulfonate}{tosylaat} en daarna behandeld met
natriumethoxide in ethanol. Geef de \omterm{def:b1:stereochemistry-in-depth:isomers}{configuratie} van de ether. Wat zou je
krijgen door dezelfde alcohol met zwavelzuur te verwarmen?
\end{exercise}

\begin{exercise}[$\star\star$]\label{exo:b1:alcohol-activation:7}
Schrijf het mechanisme op van de reactie van 2-methylpropaan-2-ol met
geconcentreerd zoutzuur en verklaar waarom de reactie bij
kamertemperatuur snel verloopt.
\end{exercise}

\begin{exercise}[$\star\star$]\label{exo:b1:alcohol-activation:8}
Schrijf het mechanisme op van de vorming van ethoxyethaan uit ethanol in
zwavelzuur, en de concurrerende eliminatie.
\end{exercise}

\begin{exercise}[$\star\star$]\label{exo:b1:alcohol-activation:9}
Welke alcohol wordt in de classificatietest van dit hoofdstuk als eerste
troebel: butaan-1-ol, butaan-2-ol of 2-methylpropaan-2-ol? Schrijf het
product op.
\end{exercise}

\begin{exercise}[$\star\star\star$]\label{exo:b1:alcohol-activation:10}
3,3-Dimethylbutaan-2-ol geeft bij verwarming in zuur milieu vooral
2,3-dimethylbut-2-een, waarvan het koolstofskelet verschilt van dat van de
alcohol. Stel een mechanisme voor met een 1,2-verschuiving van een
methylgroep in het \omterm{def:b1:electronic-effects:intermediates}{carbokation}, en verantwoord die verschuiving.
\end{exercise}

\begin{exercise}[$\star\star\star$]\label{exo:b1:alcohol-activation:11}
Ontwerp een synthese van \cip{S}-2-methoxybutaan uit \cip{R}-butaan-2-ol,
en een andere uit \cip{S}-butaan-2-ol, telkens met de \omterm{def:b1:stereochemistry-in-depth:isomers}{configuratie} onder
controle.
\end{exercise}

\begin{exercise}[$\star\star\star$]\label{exo:b1:alcohol-activation:12}
Toon met het \omterm{def:b1:extent-q-and-k:quotient}{reactiequotiënt} (\cref{ch:b1:extent-q-and-k}) aan dat het
evenwicht \ce{2CH3CH2OH <=> (CH3CH2)2O + H2O} naar de ether gedreven kan
worden door continu water te verwijderen.
\end{exercise}

\section{Probleem: twee wegen naar een ether}

\begin{problem}[{Weekendprobleem --- de twee williamsonroutes naar
methyl-tert-butylether, een ether met gecontroleerde stereochemie uit een
optisch actieve alcohol, de dehydratatie van een tertiaire alcohol, en de
verkregen massa ether}]
\label{pb:b1:alcohol-activation:1}
2-Methoxy-2-methylpropaan (methyl-tert-butylether, MTBE) werd op grote
schaal geproduceerd als brandstofadditief. Molaire massa's (\unit{g/mol}):
\ce{C4H10O} 74.0, \ce{C5H12O} 88.0, \ce{CH3I} 141.9.

\medskip
\noindent\textbf{Deel I --- MTBE via de \omterm{def:b1:alcohol-activation:williamson}{ethersynthese van Williamson}.}
\begin{enumerate}
  \item Schrijf de twee paren partners op (\omterm{def:b1:alcohol-activation:alkoxide}{alkoxide} en halogenide) die
        MTBE zouden kunnen geven.
  \item Welk paar faalt? Schrijf de reactie op die in plaats daarvan
        verloopt.
  \item Schrijf de geslaagde reactie en haar mechanisme op.
  \item Hoe wordt natrium-tert-butoxide bereid uit 2-methylpropaan-2-ol?
        Waarom niet met natriumhydroxide?
  \item Waarom wordt de reactie uitgevoerd in de droge alcohol of in een
        \omterm{def:b1:intermolecular-forces-solvents:solvent-classes}{aprotisch oplosmiddel}?
  \item Is het product \omterm{def:b1:stereochemistry-in-depth:stereocentre}{chiraal}?
\end{enumerate}

\medskip
\noindent\textbf{Deel II --- Een optisch actieve ether.}
\begin{enumerate}[resume]
  \item \cip{R}-Butaan-2-ol wordt omgezet in zijn \omterm{def:b1:alcohol-activation:sulfonate}{tosylaat}. Schrijf de
        reactie op en geef de \omterm{def:b1:stereochemistry-in-depth:isomers}{configuratie} van het \omterm{def:b1:alcohol-activation:sulfonate}{tosylaat}.
  \item Het \omterm{def:b1:alcohol-activation:sulfonate}{tosylaat} reageert met natriummethoxide. Schrijf het product
        en zijn \omterm{def:b1:stereochemistry-in-depth:isomers}{configuratie} op.
  \item Welke \omterm{def:b1:stereochemistry-in-depth:isomers}{configuratie} van de alcohol geeft langs dezelfde weg
        \cip{R}-2-methoxybutaan?
  \item Een andere weg: het \omterm{def:b1:alcohol-activation:alkoxide}{alkoxide} van \cip{R}-butaan-2-ol met
        joodmethaan. \omterm{def:b1:stereochemistry-in-depth:isomers}{Configuratie} van het product? Waarom?
  \item Welke nevenreactie concurreert in vraag 8, en waarom blijft ze
        hier beperkt?
  \item Waarom zou het verwarmen van \cip{R}-butaan-2-ol met methanol in
        zwavelzuur een bijna racemische ether geven, als er al ether
        ontstaat?
\end{enumerate}

\medskip
\noindent\textbf{Deel III --- \omterm{def:b1:alcohol-activation:dehydration}{Dehydratatie} van een tertiaire alcohol.}
\begin{enumerate}[resume]
  \item Schrijf het mechanisme op van de \omterm{def:b1:alcohol-activation:dehydration}{dehydratatie} in zuur milieu van
        2-methylbutaan-2-ol.
  \item Geef de twee alkenen en voorspel het hoofdproduct.
  \item Waarom ontstaat uit deze alcohol geen ether?
  \item Waarom wordt 2-methylbut-2-een afgedestilleerd naarmate het
        ontstaat?
  \item Teken het \omterm{def:b1:elementary-steps:energy-profile}{energieprofiel} van de E1-stap en duid de
        \omterm{def:b1:elementary-steps:rds}{snelheidsbepalende stap} aan.
  \item Zou een primaire alcohol volgens hetzelfde mechanisme
        gedehydrateerd kunnen worden?
\end{enumerate}

\medskip
\noindent\textbf{Deel IV --- Massa's.}
\begin{enumerate}[resume]
  \item \qty{10.0}{g} 2-methylpropaan-2-ol wordt omgezet in het \omterm{def:b1:alcohol-activation:alkoxide}{alkoxide}.
        Bereken de hoeveelheid stof.
  \item Welke massa joodmethaan is nodig voor een overmaat van
        \qty{10}{\percent}?
  \item Bereken de theoretische massa MTBE.
  \item Bereken de verkregen massa bij een \omterm{def:b1:extent-q-and-k:fractional-extent}{rendement} van
        \qty{75}{\percent}.
\end{enumerate}
\end{problem}
